\medskip

Dans tout l'exercice, les probabilités seront, si nécessaire, arrondies à $10^{-3}$ près.

Une donnée binaire est une donnée qui ne peut prendre que deux valeurs:  0 ou 1.

Une donnée de ce type est transmise successivement d'une machine à une autre.

Chaque machine transmet la donnée reçue soit de manière fidèle, c'est-à-dire en
transmettant l'information telle qu'elle l'a reçue (1 devient 1 et 0 devient 0), soit de façon
contraire (1 devient 0 et 0 et devient 1).

La transmission est fidèle dans 90\,\% des cas, et donc contraire dans 10\,\% des cas.

Dans tout l'exercice, la première machine reçoit toujours la valeur 1.

\medskip

\textbf{Partie A}

\medskip

Pour tout entier naturel $n \geqslant 1$, on note :

\begin{itemize}[label=$\bullet~$]
\item $V_n$ l'évènement : \og la $n$-ième machine détient la valeur 1 \fg{};
\item $\overline{V_n}$ l'évènement : \og la $n$-ième machine détient la valeur 0 \fg{}.
\end{itemize}

\medskip

\begin{enumerate}
\item 
	\begin{enumerate}
		\item Recopier et compléter l'arbre de probabilité ci-dessous.

\medskip

\begin{center}
\pstree[treemode=R,levelsep=3cm]{\TR{$V_1$~}}
{\pstree{\TR{$V_2~$}\taput{\ldots}}
	{\TR{$V_3$} \taput{\ldots}
	\TR{$\overline{V_3}$} \tbput{\ldots}
	}
\pstree{\TR{$\overline{V_2}~$}\tbput{\ldots}}
	{\TR{$V_3$} \taput{\ldots}
	\TR{$\overline{V_3}$} \tbput{\ldots}
	}
}
\end{center}

		\item Démontrer que $P(V_3) = 0,82$ et interpréter ce résultat dans le contexte de l'exercice.
		\item Sachant que la troisième machine a reçu la valeur 1, calculer la probabilité que la deuxième machine ait aussi reçu la valeur 1.
	\end{enumerate}
\item Pour tout entier naturel $n \geqslant 1$, on note $p_n = P(V_n)$.

La première machine a reçu la valeur 1, on a donc $p_1 = 1$.
	\begin{enumerate}
		\item Démontrer que pour tout entier naturel $n \geqslant 1$ :
		
\[p_{n+1} = 0,8p_n + 0,1.\]

		\item Démontrer par récurrence que pour tout entier naturel $n \geqslant 1$,

\[p_n = 0,5 \times 0,8^{n-1} + 0,5.\]

		\item Calculer la limite de $p_n$ lorsque $n$ tend vers l'infini. Interpréter ce résultat dans le contexte de l'exercice.
	\end{enumerate}
\end{enumerate}


\textbf{Partie B}

\medskip

Pour modéliser en langage Python la transmission de la donnée binaire décrite en début d'exercice, on considère la fonction \texttt{simulation} qui prend en paramètre un entier naturel $n$ qui représente le nombre de transmissions réalisées d'une machine à une autre, et qui renvoie la liste des valeurs successives de la donnée binaire.

On donne ci-dessous le script incomplet de cette fonction.

On rappelle que l'instruction \texttt{rand()} renvoie un nombre aléatoire de l'intervalle [0~;~1[.


\[
\begin{array}{l}
\texttt{def simulation(n):}\\
\texttt{\hspace{1em}donnee = 1}\\
\texttt{\hspace{1em}liste = [donnee]}\\
\texttt{\hspace{1em}for k in range(n):}\\
\texttt{\hspace{2em}if rand() < 0.1:}\\
\texttt{\hspace{3em}donnee = 1 - donnee}\\
\texttt{\hspace{2em}liste.append(donnee)}\\
\texttt{\hspace{1em}return liste}
\end{array}
\]

\medskip

Par exemple, \texttt{simulation(3)} peut renvoyer [1,\: 0,\: 0,\: 1].Cette liste traduit :

\begin{itemize}[label=$\bullet~$]
\item qu'une donnée binaire a été successivement transmise trois fois entre quatre machines;
\item la première machine qui détient la valeur 1 a transmis de façon contraire cette donnée à la deuxième machine ;
\item la deuxième machine a transmis la donnée qu'elle détient de façon fidèle à la troisième ;
\item la troisième machine a transmis de façon contraire la donnée qu'elle détient à la quatrième.
\end{itemize}

\medskip

\begin{enumerate}
\item Déterminer le rôle des instructions des lignes 5 et 6 de l'algorithme ci-dessus.
\item Calculer la probabilité que simulation(4) renvoie la liste [1,\:1,\:1,\:1,\:1] et la probabilité que \texttt{simulation(6)} renvoie la liste [1,\: 0,\:1,\:0,\:0,\:1,\:1].
\end{enumerate}

\bigskip
